\originalchapter{分析与计算的联系}\label{chap:13}
本章对应原第25讲复习与收束, 第317--339页. 该讲把全课程重新归为凸分析、对偶、算法三部分. 以下以可核的解证书串联它们, 不重复抄写前文已经证明的结果.

\originalsection{从模型到证书}
一个凸问题能否求解, 至少有四个不同问题. 首先, 是否存在可行点并且最优值有限; 其次, 最优值是否取到; 然后, 是否有零对偶间隙与最优乘子; 最后, 计算过程是否以可证明的速度产生足够好的近似. 凸性只使局部最优等于全局最优, 并不自动回答后三个问题.

第\ref{chap:1}章把凸性变成全局切平面下界与投影条件. 第\ref{chap:2}章用相对内部控制低维域, 用衰退锥判定紧水平集与解的存在. 第\ref{chap:3}章说明线性像和部分极小化可能丢失闭性, 并以分离与共轭把点集描述转换成半空间描述. 第\ref{chap:4}章的扰动函数将零对偶间隙解释为原点的高度不在闭凸化中下降, 将乘子存在解释为原点有非竖直支持平面. 第\ref{chap:5}章把这个支持斜率写成次梯度, 将其用于最优性与灵敏度.

对可行候选 $x$ 与对偶可行乘子, 原目标是上界, 对偶值是下界. 二者之差才是直接的停止证书. 若有不可行迭代, 目标值可能低于原最优值, 不能独立当成上界. 次梯度距离估计、近端势函数和障碍的原始--对偶间隙, 分别提供不同的可证明进展量. “目标看起来下降”与“已达到指定精度”不是同一个结论.

\originalsection{算法利用的结构}
当目标不可微而次梯度廉价, 第\ref{chap:7}章以距离递推分析投影次梯度. 它允许目标上升, 步长选择影响值误差与点收敛, 所以逐步下降并非必要的进展指标. 当支持平面容易保存, 第\ref{chap:8}章建立外部多面体下模型; 当线性优化子问题容易, 单纯形分解维护内部可行凸包. 两者在共轭下相互对应, 但它们给出的界方向和子问题条件不同.

第\ref{chap:9}章通过二次正则化产生唯一的近端点, 并用到最优集合的距离解释其收敛. 在对偶侧应用同一近端步骤, 就出现增广拉格朗日乘子更新; 加上多面体近似则产生束方法. 第\ref{chap:10}章的障碍法从严格可行域内部逼近边界, 其参数趋零可能带来病态, 需要适合的内层求解. 大和目标则可利用循环或随机增量, 此时单次工作量减少, 但固定步长可能只到误差邻域.

对光滑目标, 第\ref{chap:11}章的二次上界使梯度投影由次梯度的 $O(\varepsilon^{-2})$ 次数改进到 $O(\varepsilon^{-1})$; 外推与势函数进一步给 $O(\varepsilon^{-1/2})$ 次数阶. 这些次数界以相应梯度、投影或近端查询为单位, 并不忽略每次查询的成本. 非光滑平滑化需要同时计算近似误差与平滑函数的Lipschitz常数. Bregman正则化和熵镜像下降适应单纯形等几何结构, 但不能不核对定义域、正分量与三点不等式条件就直接搬用欧氏分析.

线性规划还具有有限基结构, 第\ref{chap:12}章的单纯形法沿基换元, 用约化成本与比值检验给最优或无界证书. 退化时算法可能换基而不移动点, 防循环规则在这里有实际数学作用. 灵敏度则是在基保持可行与最优的范围内, 把对偶乘子变为精确的参数斜率.

\originalsection{一个复合模型}
\begin{example}
对 $\lambda>0$, 考虑 $F(x)=\tfrac12\norm{Ax-b}^2+\lambda\norm{x}_1$. 给出最优性证书、一个可实现迭代及可证明的停止判断.
\end{example}
\begin{solution}
光滑部分梯度为 $A\T(Ax-b)$, Lipschitz常数可取 $L=\norm A^2$. 全空间有限连续与闭真 $\ell_1$ 项满足和法则, 所以最优性等价于存在 $s\in\partial\norm{x}_1$ 使 $A\T(Ax-b)+\lambda s=0$. 各坐标上, $x_j\ne0$ 时 $s_j=\operatorname{sgn}(x_j)$, $x_j=0$ 时 $s_j\in[-1,1]$. 因而非零坐标满足精确梯度等式, 零坐标满足 $|[A\T(Ax-b)]_j|\le\lambda$.

取 $0<\alpha\le1/L$ (若 $A=0$, 直接 $x=0$ 最优), 迭代 $z=x-\alpha A\T(Ax-b)$, 再逐坐标取 $x_j^+=\operatorname{sgn}(z_j)(|z_j|-\alpha\lambda)_+$. 这是梯度近端法, 软阈值由一维次梯度条件直接求得. 若使用外推, 应按第\ref{chap:11}章的参数与势函数证明分析, 不能只在该式前随意加一个动量项.

还可构造一个直接的对偶界. 将残差变量 $r=Ax-b$ 拆出并用乘子 $y$, 极小化 $\tfrac12\norm r^2+\lambda\norm x_1+y\T(Ax-b-r)$, 得对偶 $q(y)=-\tfrac12\norm y^2-b\T y$, 约束 $\norm{A\T y}_\infty\le\lambda$. 对任意残差候选 $r$, 若 $\norm{A\T r}_\infty>\lambda$, 将其按比例缩成对偶可行 $y$; 若已满足, 直接取 $y=r$. 因而 $F(x)-q(y)\ge0$ 是一个可计算误差上界. 仅检查软阈值后坐标变动小, 并不等于检查了这一目标误差.
\end{solution}
这个例子同时使用凸性、次微分、对偶、近端与次数分析. 它说明全书各部分的联系, 也说明模型与算法的每项假设都应在实际应用处重新核对.
